\documentclass[11pt,a4paper]{article}

\usepackage[margin=1in]{geometry}
\usepackage{booktabs}
\usepackage{longtable}
\usepackage{array}
\usepackage{listings}
\usepackage{xcolor}
\usepackage{hyperref}
\usepackage{enumitem}
\usepackage{amsmath}
\usepackage{amssymb}
\usepackage[T1]{fontenc}
\usepackage{lmodern}

\hypersetup{
    colorlinks=true,
    linkcolor=blue,
    urlcolor=blue,
    citecolor=blue
}

\lstset{
    basicstyle=\ttfamily\small,
    breaklines=true,
    frame=single,
    columns=fullflexible,
    keepspaces=true
}

\title{\textbf{Self-Consistency vs. Source-Grounded NLI for Hallucination Detection}}
\author{}
\date{}

\begin{document}

\maketitle

\begin{abstract}
This project presents an empirical comparison of two black-box hallucination-detection strategies for LLM-generated summaries. Method A uses behavioral self-consistency checking, while Method B performs direct source-grounded natural language inference (NLI). The benchmark evaluates where self-consistency succeeds and fails relative to direct textual evidence.
\end{abstract}

\section{Key Findings}

\subsection{Research Question}

For claims injected into LLM-generated summaries:

\begin{quote}
How does the catch rate of black-box self-consistency checking
(\textbf{Method A}) vary across claim types, and does source-grounded
entailment (\textbf{Method B}) catch a distinct set of failures?
\end{quote}

\textbf{Model under test:} \texttt{openai/gpt-4o-mini}, accessed through the OpenRouter API for both generation and verification.

\subsection{Catch Rate (Recall) on Injected Unsupported Claims}

\begin{table}[h]
\centering
\begin{tabular}{lccc}
\toprule
\textbf{Category} & \textbf{N} & \textbf{Method A} & \textbf{Method B} \\
\midrule
Numerical fabrication & 11 & 100.0\% (11/11) & 100.0\% (11/11) \\
Entity fabrication & 10 & 90.0\% (9/10) & 100.0\% (10/10) \\
Causal overreach & 25 & 96.0\% (24/25) & 100.0\% (25/25) \\
\midrule
\textbf{Overall} & \textbf{46} & \textbf{95.7\% (44/46)} & \textbf{100.0\% (46/46)} \\
\bottomrule
\end{tabular}
\caption{Recall on injected unsupported claims.}
\end{table}

\subsection{False Positive Rate on the 86 Supported Controls}

\begin{table}[h]
\centering
\begin{tabular}{lcc}
\toprule
\textbf{Method} & \textbf{False Positives} & \textbf{FPR} \\
\midrule
Method A: Self-Consistency & 5/86 & 5.8\% \\
Method B: Source-Grounded NLI & 4/86 & 4.7\% \\
\bottomrule
\end{tabular}
\caption{False positive rate on supported control claims.}
\end{table}

\subsection{Cross-Tabulation: Unsupported Claims Only}

\begin{table}[h]
\centering
\begin{tabular}{lcc}
\toprule
 & \textbf{Method B Caught} & \textbf{Method B Missed} \\
\midrule
\textbf{Method A Caught} & 44 & 0 \\
\textbf{Method A Missed} & 2 & 0 \\
\bottomrule
\end{tabular}
\caption{Cross-tabulation of unsupported claims.}
\end{table}

\subsection{Takeaway}

Method B strictly dominates on recall: it catches \textbf{100\% of injected unsupported claims}, including the two claims missed by Method A.

Method A, despite having no access to the source document during verification, achieves \textbf{95.7\% recall}. Its failures occur primarily when an LLM reproduces a plausible narrative schema across independently generated summaries and the equivalence verifier mistakes this consistency for factual corroboration.

The results highlight a fundamental distinction:

\begin{quote}
\textbf{Self-consistency measures consensus; source-grounded NLI measures textual support.}
\end{quote}


\section{Repository Architecture}

\begin{lstlisting}
.
├── .env.example
├── .gitignore
├── README.md
├── research_report.md
│
├── claims_dataset.py
├── claims_dataset.json
│
├── method_a_self_consistency.py
├── method_b_entailment.py
│
├── evaluate_and_score.py
├── check_missed_claims.py
│
├── regenerations_cache.json
├── method_a_results.json
├── method_b_results.json
└── evaluation_summary.json
\end{lstlisting}

\begin{itemize}
    \item \texttt{.env.example}: Template for environment variables and API keys.
    \item \texttt{.gitignore}: Ignores \texttt{.env}, \texttt{.venv}, caches, and IDE configurations.
    \item \texttt{README.md}: Project overview, methodology, results, and reproduction guide.
    \item \texttt{research\_report.md}: Detailed research write-up and error analysis.
    \item \texttt{claims\_dataset.py}: Defines the 132-claim adversarial benchmark.
    \item \texttt{method\_a\_self\_consistency.py}: Implements Method A.
    \item \texttt{method\_b\_entailment.py}: Implements Method B.
    \item \texttt{evaluate\_and\_score.py}: Calculates evaluation metrics.
    \item \texttt{check\_missed\_claims.py}: Inspects Method A false negatives.
\end{itemize}

The files \texttt{.env}, \texttt{.venv/}, and \texttt{\_\_pycache\_\_/} are deliberately excluded from the repository through \texttt{.gitignore}.

\subsection{Dataset Generation}

The benchmark JSON file is generated by running:

\begin{lstlisting}[language=bash]
python claims_dataset.py
\end{lstlisting}

The downstream evaluation scripts read the serialized JSON directly from disk.


\section{Methodology}

\subsection{Adversarial Dataset}

The benchmark is defined in \texttt{claims\_dataset.py} and contains:

\begin{itemize}
    \item \textbf{132 total claims}
    \item \textbf{13 source documents}
    \item \textbf{46 unsupported injected claims}
    \item \textbf{86 supported control claims}
\end{itemize}

The source documents cover multiple technical and factual domains, including:

\begin{itemize}
    \item Astronomy press releases
    \item Corporate earnings reports
    \item Municipal continuity orders
    \item Clinical trial readouts
    \item Semiconductor hardware announcements
\end{itemize}

Each dataset entry contains:

\begin{lstlisting}
claim_id
doc_id
source_text
claim_text
category
is_unsupported
\end{lstlisting}

\subsection{Claim Categories}

\subsubsection{Unsupported Injections -- N=46}

\begin{table}[h]
\centering
\begin{tabular}{lcl}
\toprule
\textbf{Category} & \textbf{N} & \textbf{Description} \\
\midrule
\texttt{numerical} & 11 & Altered metrics, dates, percentages, and sample sizes \\
\texttt{entity} & 10 & Fabricated researchers, organizations, partners, or instruments \\
\texttt{causal} & 25 & Unsupported causal extrapolations and premature efficacy claims \\
\bottomrule
\end{tabular}
\caption{Unsupported claim categories.}
\end{table}

\subsubsection{Supported Controls -- N=86}

Supported claims are grounded factual claims and appropriately hedged statements where:

\begin{lstlisting}
is_unsupported = False
\end{lstlisting}


\subsection{Method A -- Self-Consistency}

\textbf{Implementation:}

\begin{lstlisting}
method_a_self_consistency.py
\end{lstlisting}

Method A is deliberately \textbf{source-blind at verification time}. It attempts to detect hallucinations by measuring whether an injected claim consistently appears across independently generated summaries.

\subsubsection{Procedure}

\begin{enumerate}
    \item For each of the 13 source documents, generate \textbf{5 alternative summaries}.
    \item Generation uses:
\end{enumerate}

\begin{lstlisting}
N_REGENERATIONS = 5
temperature = 0.7
\end{lstlisting}

\begin{enumerate}[resume]
    \item Generated summaries are cached in \texttt{regenerations\_cache.json}.
    \item For each of the 132 claims, run a binary equivalence check against each of the 5 summaries.
    \item Each equivalence check uses:
\end{enumerate}

\begin{lstlisting}
temperature = 0.0
\end{lstlisting}

\begin{enumerate}[resume]
    \item Calculate the agreement score:
\end{enumerate}

\begin{lstlisting}
agreement_score = matches / 5
\end{lstlisting}

\begin{enumerate}[resume]
    \item Flag the claim as unsupported when:
\end{enumerate}

\begin{lstlisting}
agreement_score <= 0.40
\end{lstlisting}

In other words, a claim is considered ``caught'' when it appears to be supported by no more than \textbf{2 of the 5 generated summaries}.

\subsubsection{Important Property}

Method A does \textbf{not} inspect the original source text during verification. Its signal is therefore behavioral consistency rather than direct factual grounding.


\subsection{Method B -- Source-Grounded NLI}

\textbf{Implementation:}

\begin{lstlisting}
method_b_entailment.py
\end{lstlisting}

Method B performs a direct source-grounded entailment check. For each claim, the verifier receives the relevant source text and determines whether the source directly supports the claim.

\subsubsection{Procedure}

For each claim:

\begin{enumerate}
    \item Provide the source text.
    \item Provide the claim.
    \item Ask whether the source directly supports the assertion.
    \item Run the verification prompt at:
\end{enumerate}

\begin{lstlisting}
temperature = 0.0
\end{lstlisting}

\begin{enumerate}[resume]
    \item Parse the categorical output:
\end{enumerate}

\begin{lstlisting}
VERDICT: YES
\end{lstlisting}

or

\begin{lstlisting}
VERDICT: NO
\end{lstlisting}

\begin{enumerate}[resume]
    \item A \texttt{NO} verdict marks the unsupported claim as caught.
\end{enumerate}

Method B therefore has direct access to the evidence against which the claim is evaluated.


\subsection{Scoring}

Evaluation is performed by:

\begin{lstlisting}
evaluate_and_score.py
\end{lstlisting}

The script joins:

\begin{lstlisting}
method_a_results.json
method_b_results.json
\end{lstlisting}

using:

\begin{lstlisting}
claim_id
\end{lstlisting}

It then calculates:

\begin{itemize}
    \item Overall recall
    \item Category-level recall
    \item False positive rate on supported controls
    \item Cross-tabulation between Method A and Method B
\end{itemize}

Qualitative inspection of Method A's false negatives is performed by:

\begin{lstlisting}
check_missed_claims.py
\end{lstlisting}


\section{Detailed Results and Failure Analysis}

Method A missed exactly \textbf{2 of the 46 injected unsupported claims}. Both were correctly caught by Method B.

\subsection{Failure Case 1 -- Entity Fabrication}

\textbf{Claim ID:} \texttt{doc11\_c04}

\textbf{Category:} Entity Fabrication

\textbf{Method A Agreement Score:} 0.8 / 80\%

\subsubsection{Injected Claim}

\begin{quote}
``City Attorney Jennifer Martinez advised the council that state law allows local age restrictions on e-bike operators.''
\end{quote}

\subsubsection{Source Truth}

The source states that councilmembers asked municipal staff to pursue legislative options. No city attorney named Jennifer Martinez was identified as having advised the council.

\subsubsection{Why Method A Failed}

The generated summaries repeatedly focused on the municipal regulatory debate. The equivalence verifier therefore encountered a plausible narrative structure involving:

\begin{itemize}
    \item Municipal regulation
    \item City officials
    \item State law
    \item E-bike restrictions
\end{itemize}

Because these concepts appeared consistently across generations, the verifier treated the fabricated official attribution as semantically supported.

\subsubsection{Failure Mechanism}

The verifier exhibited \textbf{semantic drift}:

\begin{lstlisting}
Topical consistency
        |
        v
Similar narrative framing
        |
        v
Perceived semantic equivalence
        |
        v
False acceptance
\end{lstlisting}

The problem is that contextual plausibility does not establish factual support.


\subsection{Failure Case 2 -- Causal Overreach}

\textbf{Claim ID:} \texttt{doc13\_c06}

\textbf{Category:} Causal Overreach

\textbf{Method A Agreement Score:} 1.0 / 100\%

\subsubsection{Injected Claim}

\begin{quote}
``Webb's coronagraph technology enabled this discovery, which would have been completely impossible with any other existing telescope.''
\end{quote}

\subsubsection{Source Truth}

The source describes how the coronagraph operates and its role in the observation. However, it does not establish that:

\begin{itemize}
    \item the discovery was exclusively enabled by Webb's coronagraph, or
    \item the discovery would have been completely impossible using every other existing telescope.
\end{itemize}

\subsubsection{Why Method A Failed}

Every generated summary emphasized the importance of the coronagraph. The equivalence verifier therefore observed extremely high agreement:

\begin{lstlisting}
5 / 5 summaries
=
agreement score of 1.0
\end{lstlisting}

However, the injected claim contained a much stronger assertion than the source:

\begin{lstlisting}
"The coronagraph was useful"
\end{lstlisting}

became:

\begin{lstlisting}
"The coronagraph made the discovery uniquely possible."
\end{lstlisting}

\subsubsection{Failure Mechanism}

Method A conflated repeated rhetorical emphasis with evidence for an absolute causal claim.

This is particularly problematic for claims containing:

\begin{itemize}
    \item ``only''
    \item ``completely impossible''
    \item ``the first''
    \item ``the only''
    \item ``always''
    \item ``never''
    \item ``caused''
    \item ``enabled''
    \item ``proved''
\end{itemize}


\section{What the Failure Cases Demonstrate}

The two false negatives expose the central theoretical limitation of self-consistency-based hallucination detection:

\begin{quote}
\textbf{A model can consistently repeat a hallucination.}
\end{quote}

Independent generations are not independent sources of truth. They are samples from the same model's learned distribution.

If the model has a strong prior about a particular narrative structure, multiple generations can converge on the same unsupported assertion.

Therefore:

\[
\text{High agreement} \neq \text{High factuality}
\]

More precisely:

\[
\begin{array}{c}
\text{Self-consistency}\\
\downarrow\\
\text{Behavioral agreement}
\end{array}
\qquad
\begin{array}{c}
\text{Source-grounded NLI}\\
\downarrow\\
\text{Evidence support}
\end{array}
\]

This distinction explains why Method B was able to catch both of Method A's false negatives.


\section{Quickstart}

\subsection{Requirements}

\begin{itemize}
    \item Python 3.10 or higher
    \item An OpenRouter API key (or OpenAI-compatible provider)
\end{itemize}

\subsection{Clone the Repository}

\begin{lstlisting}[language=bash]
git clone <YOUR_REPOSITORY_URL>
cd <YOUR_REPOSITORY_DIRECTORY>
\end{lstlisting}

\subsection{Create and Activate a Virtual Environment}

\subsubsection{macOS / Linux}

\begin{lstlisting}[language=bash]
python3 -m venv .venv
source .venv/bin/activate
\end{lstlisting}

\subsubsection{Windows}

\begin{lstlisting}[language=bash]
python -m venv .venv
.venv\Scripts\activate
\end{lstlisting}

\subsection{Install Dependencies}

\begin{lstlisting}[language=bash]
pip install openai tqdm python-dotenv
\end{lstlisting}

\subsection{Configure API Credentials}

Copy the example environment file:

\begin{lstlisting}[language=bash]
cp .env.example .env
\end{lstlisting}

Then edit \texttt{.env}:

\begin{lstlisting}
OPENROUTER_API_KEY=your_key_here
\end{lstlisting}

Do \textbf{not} commit \texttt{.env} to GitHub. The repository's \texttt{.gitignore} is configured to exclude it.


\section{Reproduction Pipeline}

Run the scripts sequentially.

\subsection{Step 1 -- Generate Dataset}

\begin{lstlisting}[language=bash]
python claims_dataset.py
\end{lstlisting}

This creates:

\begin{lstlisting}
claims_dataset.json
\end{lstlisting}

\subsection{Step 2 -- Run Method A}

\begin{lstlisting}[language=bash]
python method_a_self_consistency.py
\end{lstlisting}

This generates:

\begin{lstlisting}
regenerations_cache.json
method_a_results.json
\end{lstlisting}

The first file caches the five generated summaries per source document.

\subsection{Step 3 -- Run Method B}

\begin{lstlisting}[language=bash]
python method_b_entailment.py
\end{lstlisting}

This generates:

\begin{lstlisting}
method_b_results.json
\end{lstlisting}

\subsection{Step 4 -- Evaluate Both Methods}

\begin{lstlisting}[language=bash]
python evaluate_and_score.py
\end{lstlisting}

This generates:

\begin{lstlisting}
evaluation_summary.json
\end{lstlisting}

The evaluation contains:

\begin{itemize}
    \item Category-level recall
    \item Overall recall
    \item False positive rates
    \item Cross-tabulation
\end{itemize}

\subsection{Step 5 -- Inspect Method A Misses}

\begin{lstlisting}[language=bash]
python check_missed_claims.py
\end{lstlisting}

This prints the unsupported claims that:

\begin{itemize}
    \item Method A missed
    \item Method B caught
\end{itemize}

This provides a qualitative error analysis of the self-consistency approach.


\section{Caching}

\texttt{method\_a\_self\_consistency.py} caches generated summaries in:

\begin{lstlisting}
regenerations_cache.json
\end{lstlisting}

This prevents unnecessary regeneration when rerunning the evaluation.

To regenerate all summaries from scratch, delete the cache:

\begin{lstlisting}[language=bash]
rm regenerations_cache.json
\end{lstlisting}

Then rerun:

\begin{lstlisting}[language=bash]
python method_a_self_consistency.py
\end{lstlisting}


\section{Environment Variables}

The project expects:

\begin{lstlisting}
OPENROUTER_API_KEY=your_key_here
\end{lstlisting}

The scripts use \texttt{python-dotenv} to load the value from \texttt{.env}.

Example \texttt{.env.example}:

\begin{lstlisting}
OPENROUTER_API_KEY=
\end{lstlisting}


\section{Reproducibility Notes}

The benchmark is deterministic at the verification stage because both verification methods use:

\begin{lstlisting}
temperature = 0.0
\end{lstlisting}

Method A's summary generation intentionally uses stochastic generation:

\begin{lstlisting}
temperature = 0.7
N = 5
\end{lstlisting}

Therefore, deleting \texttt{regenerations\_cache.json} and regenerating summaries can potentially produce different intermediate generations and agreement scores.

For reproducing the reported results as closely as possible, use the committed:

\begin{lstlisting}
regenerations_cache.json
method_a_results.json
method_b_results.json
evaluation_summary.json
\end{lstlisting}


\section{Threats to Validity}

\subsection{Single-Model Evaluator}

Both summary generation and verification rely on:

\begin{lstlisting}
openai/gpt-4o-mini
\end{lstlisting}

Therefore, the results do not establish that the same performance would occur with other models. Cross-model verification using models such as Claude, Gemini, or Llama was not evaluated.

\subsection{Dataset Scale}

The benchmark contains:

\begin{lstlisting}
132 claims
13 source documents
46 unsupported claims
86 supported controls
\end{lstlisting}

The dataset was constructed by a single author. No multi-annotator inter-rater reliability analysis, such as Cohen's Kappa, was performed.

\subsection{Threshold Parameterization}

Method A uses:

\begin{lstlisting}
N = 5 regenerations
threshold <= 0.40
\end{lstlisting}

This represents an explicit design choice. Changing either parameter could alter the precision/recall trade-off.

For example:

\begin{lstlisting}
More regenerations
        |
        v
Higher computational cost
        |
        v
Potentially more stable agreement estimates
\end{lstlisting}

Changing the threshold can similarly alter how aggressively low-consensus claims are flagged.

\subsection{Prompt Dependence}

Both methods depend on the exact prompts used for:

\begin{itemize}
    \item Summary generation
    \item Equivalence verification
    \item Source-grounded entailment
\end{itemize}

Small changes to prompt wording may affect model behavior.

\subsection{Synthetic Adversarial Claims}

The benchmark consists of deliberately injected unsupported claims. Although the categories are designed to represent common hallucination patterns, performance on naturally occurring hallucinations in production systems may differ.


\section{Interpretation}

The benchmark suggests that source-grounded verification provides a stronger signal for detecting unsupported claims than self-consistency alone.

\subsection{Method A}

\textbf{Advantages:}

\begin{itemize}
    \item Does not require source access during verification
    \item Can detect claims that fail to recur across generations
    \item Achieves high recall despite being source-blind
    \item Can potentially be applied when only generated outputs are available
\end{itemize}

\textbf{Limitations:}

\begin{itemize}
    \item Consensus can reinforce the same hallucination
    \item Vulnerable to plausible entity fabrication
    \item Vulnerable to causal overreach
    \item Requires multiple generations
    \item Higher inference cost due to repeated generation and verification
\end{itemize}

\subsection{Method B}

\textbf{Advantages:}

\begin{itemize}
    \item Directly evaluates claims against source evidence
    \item Achieved 100\% recall on this benchmark
    \item Caught both Method A false negatives
    \item Requires only a single verification pass per claim
\end{itemize}

\textbf{Limitations:}

\begin{itemize}
    \item Requires access to the source document
    \item Depends on the verifier's ability to interpret the source correctly
    \item Still produced false positives on 4 supported controls
    \item Remains dependent on model and prompt quality
\end{itemize}


\section{Conclusion}

This benchmark provides an empirical comparison between two fundamentally different hallucination-detection strategies.

\subsection*{Method A -- Self-Consistency}

\begin{lstlisting}
Claim
  |
  v
Compare against multiple generated summaries
  |
  v
Measure recurrence
  |
  v
Flag low-consensus claims
\end{lstlisting}

\subsection*{Method B -- Source-Grounded NLI}

\begin{lstlisting}
Claim
  |
  v
Compare directly against source
  |
  v
Determine entailment
  |
  v
Flag unsupported claims
\end{lstlisting}

On the 46 injected unsupported claims:

\begin{lstlisting}
Method A: 95.7% recall
Method B: 100.0% recall
\end{lstlisting}

On the 86 supported controls:

\begin{lstlisting}
Method A: 5.8% false positive rate
Method B: 4.7% false positive rate
\end{lstlisting}

The two Method A failures demonstrate that \textbf{consistency is not equivalent to truth}. An LLM can repeatedly produce the same unsupported entity attribution or causal interpretation because the claim fits a plausible narrative pattern.

The results therefore support a practical distinction:

\begin{quote}
\textbf{Self-consistency can be a useful hallucination signal, but source-grounded verification provides a stronger basis for determining whether a claim is actually supported by evidence.}
\end{quote}


\section{Project Structure at a Glance}

\begin{lstlisting}
                 +----------------------+
                 |   Source Documents   |
                 +----------+-----------+
                            |
                            v
                 +----------------------+
                 |  Adversarial Claims  |
                 |       Dataset        |
                 +----------+-----------+
                            |
              +-------------+-------------+
              |                           |
              v                           v
    +-------------------+       +--------------------+
    |     Method A      |       |      Method B      |
    | Self-Consistency  |       | Source-Grounded NLI|
    +---------+---------+       +----------+---------+
              |                            |
              v                            v
    +-------------------+       +--------------------+
    | method_a_results  |       | method_b_results   |
    |       .json       |       |       .json        |
    +---------+---------+       +----------+---------+
              |                            |
              +-------------+--------------+
                            |
                            v
                 +----------------------+
                 |   Evaluation &       |
                 |   Error Analysis     |
                 +----------+-----------+
                            |
                            v
                 +----------------------+
                 | evaluation_summary   |
                 |        .json         |
                 +----------------------+
\end{lstlisting}


\section*{License}

MIT License. See \texttt{LICENSE} for details.

\section*{Citation}

If you use this benchmark or methodology in research, please cite this repository.

\end{document}